\documentclass[12pt,a4paper]{article}
        \makeatletter
        \def\input@path{{../}}
        \makeatother
        \usepackage{almeria-fichas}

        \FichaMeta{Sesion 08}{Area de zonas urbanas}{Bloque 2 · Geometria urbana}{Calcular el area de figuras simples y compuestas para estimar el espacio interior que ocupan.}{Procedimientos}{Comprender, Aplicar, Evaluar}

        \begin{document}

        \FichaCabecera

        \begin{materialesbox}
        \begin{itemize}
    \item Ficha impresa
    \item Lapiz
    \item Regla
    \item Calculadora sencilla si se necesita
\end{itemize}
        \end{materialesbox}

        \begin{pasosbox}
        \begin{enumerate}
    \item Lee bien si te piden perimetro o area.
    \item Escribe primero la formula adecuada.
    \item Sustituye despues los datos.
    \item Comprueba siempre la unidad final.
\end{enumerate}
        \end{pasosbox}

        \begin{recordatoriobox}
        \begin{itemize}
    \item Rectangulo: $A=b\cdot h$.
    \item Triangulo: $A=\dfrac{b\cdot h}{2}$.
    \item El area se expresa en unidades cuadradas, por ejemplo $m^2$.
\end{itemize}
        \end{recordatoriobox}

        \begin{apoyobox}
        \begin{itemize}
    \item Dibuja la base y la altura si no aparecen claras.
    \item Si la figura es compuesta, separala en partes mas simples.
    \item Rodea la unidad $m^2$ al final.
\end{itemize}
        \end{apoyobox}

        \BloomSection{comprender}

\BloomExercise{comprender}{1}{Base y altura}{Explica que representan la base y la altura en un rectangulo y en un triangulo.}{6}

\BloomExercise{comprender}{2}{Area interior}{Escribe una frase para explicar por que el area mide la parte de dentro de una figura.}{5}

\BloomSection{aplicar}

\BloomExercise{aplicar}{1}{Rectangulo}{Calcula el area de un rectangulo de base $24\,m$ y altura $15\,m$.}{5}

\BloomExercise{aplicar}{2}{Triangulo}{Calcula el area de un triangulo de base $10\,m$ y altura $6\,m$.}{5}

\BloomSection{evaluar}

\BloomExercise{evaluar}{1}{Valora una propuesta}{Evalua una posible solucion o producto relacionado con \emph{Area de zonas urbanas}. Decide si es adecuado y justifica tu opinion con criterios matematicos claros.}{8}

\BloomExercise{evaluar}{2}{Elige la mejor opcion}{Entre varias opciones posibles para cumplir este objetivo: \emph{Calcular el area de figuras simples y compuestas para estimar el espacio interior que ocupan.}, elige la mejor y argumenta por que. Incluye al menos dos criterios de valoracion.}{8}

        \begin{evidenciabox}
        \begin{itemize}
    \item Aplicar formulas de area con sentido.
    \item Distinguir la superficie interior del borde.
    \item Resolver ejemplos sencillos de figuras compuestas.
\end{itemize}
        \end{evidenciabox}

        \end{document}
